\documentclass[aspectratio=169]{beamer}

%% ============================================================
%%  pres.tex — IIST Beamer template, complete author guide
%%  Theme: beamerthemeiist.sty (v2.1, serif math)
%% ============================================================
%%  What this file demonstrates (search for the tag):
%%    [TEXT]      text formatting: bold, italic, alert, structure
%%    [LIST]      itemize, enumerate, description, with overlays
%%    [BLOCK]     block, exampleblock, alertblock, nested
%%    [COL]       columns and minipages
%%    [FIG]       figures, subfigures, side-by-side
%%    [TAB]       booktabs tables with \iistheadrow
%%    [CODE]      verbatim and listings (needs [fragile])
%%    [MATH]      inline, display, align, cases, matrix
%%    [THM]       assumption, remark, proposition
%%    [ANIM]      \pause, \onslide, \only, \uncover, \visible, \alt
%%    [TIKZ]      inline diagram
%%    [CITE]      natbib, footnote citations, bibliography
%%    [FRAME]     frame options: plain, allowframebreaks, shrink, label
%%    [WM]        watermark helpers
%%    [BACKUP]    appendix / backup slides
%%    [NOTE]      presenter notes
%% ============================================================

%% sansmathaccent + its map file are no longer needed: theme v2.1 uses
%% serif math (newpxmath). Kept here, commented, for reference.
% \usepackage{sansmathaccent}
% \pdfmapfile{+sansmathaccent.map}

\usepackage{natbib}
\bibliographystyle{dinat}
\renewcommand{\bibfont}{\tiny}

%% Bibliography file. Populate pres.bib with your references.
% \addbibresource{pres.bib}

%% ---- Theme ---------------------------------------------------------
%% Options: sans (default), serif, classicfonts,
%%          nosectionpages, subsectionpages
\usetheme{iist}
%% e.g. \usetheme[serif]{iist}
%% e.g. \usetheme[subsectionpages]{iist}

%% ---- Metadata ------------------------------------------------------
\date{\today}
\title[Short Title]{Title}
\subtitle{Subject Code: Subject Name}
\author[Presenter Name]{Presenter Name}
\authorCustom{Presenter Name}{Student Code}
\newcommand{\emailid}{youremail@domain.com}

%% Advisors are optional — comment out any you do not need.
\advisor{Role of Guide\#1}{Guide\#1 Name}{Designation of Guide\#1}
\advisorA{Role of Guide\#2}{Guide\#2 Name}{Designation of Guide\#2}

\institute[IIST]{Indian Institute of Space Science and Technology \\
Thiruvananthapuram, Kerala, India 695547 \\}

%% ---- Logo (optional; theme defaults to logo.pdf) -------------------
% \iistlogo{figures/iist-logo.pdf}
% \setlength{\iistlogoheight}{1cm}

%% ---- Section divider pages -----------------------------------------
%% The theme already inserts a styled section page at each \section.
%% To keep that and remove this custom override, delete this block.
%% The override below is the v1 look, kept for compatibility.
\AtBeginSection[]
{
    \begin{frame}
        \fontsize{16pt}{10pt}\selectfont
        \begin{center}
            \secname
        \end{center}
    \end{frame}
}

\begin{document}

%% ============================================================
%%  Title slide
%% ============================================================
\begin{frame}[plain]
    \usebeamerfont{title}
    \titlepage
\end{frame}

%% ============================================================
%%  Outline  [WM] watermark helper demo
%% ============================================================
\iistwatermarkon[0.10]   %% faded logo behind following frames
\begin{frame}
    \frametitle{Outline}
    \tableofcontents[hideallsubsections]
\end{frame}
\iistwatermarkoff        %% turn the watermark off again

%% ============================================================
\section{Text \& Formatting}  [TEXT]
%% ============================================================
\begin{frame}{Text formatting}
    Normal text. \textbf{Bold}, \textit{italic}, \texttt{monospace},
    \textsc{small caps}, \underline{underline}, \emph{emphasised}.

    \medskip
    \alert{Alerted text} uses the alert colour.
    \structure{Structured text} uses the theme colour.

    \medskip
    Coloured text: \textcolor{IISTPrimary}{IISTPrimary},
    \textcolor{IISTExample}{IISTExample},
    \textcolor{IISTAlert}{IISTAlert}.

    \medskip
    A hyperlink: \href{https://www.iist.ac.in}{IIST website}.

    \medskip
    A quote:
    \begin{quote}
        ``Simplicity is the ultimate sophistication.'' --- attributed
        to Leonardo da Vinci.
    \end{quote}
\end{frame}

%% ============================================================
\section{Lists}  [LIST]
%% ============================================================
\subsection{Itemize}
\begin{frame}{Itemize with progressive reveal}
    %% \color<.>{black} greys unvisited items, blacks the current one.
    \begin{itemize}[<+->]\color{gray}
        \item \color<.>{black} Item 1
        \item \color<.>{black} Item 2
        \begin{itemize}[<+->]\color{gray}
            \item \color<.>{black} Subitem 2.1
            \begin{itemize}[<+->]\color{gray}
                \item \color<.>{black} Subsubitem 2.1.1
                \item \color<.>{black} Subsubitem 2.1.2
            \end{itemize}
            \item \color<.>{black} Subitem 2.2
        \end{itemize}
        \item \color<.>{black} Item 3
    \end{itemize}
\end{frame}

\subsection{Enumerate}
\begin{frame}{Enumerate with overlay specs}
    \begin{enumerate}
        \item<1->      Item 1 (from slide 1)
        \item<2-3>     Item 2 (slides 2--3)
        \item<3>       Item 3 (slide 3 only)
    \end{enumerate}

    \medskip
    Plain enumerate:
    \begin{enumerate}
        \item First
        \item Second
    \end{enumerate}
\end{frame}

\subsection{Description}
\begin{frame}{Description}
    \begin{description}
        \item<1> [One] Item 1
        \item<2> [Two] Item 2
        \item<3-> [Three] Item 3
    \end{description}
\end{frame}

%% ============================================================
\section{Blocks}  [BLOCK]
%% ============================================================
\begin{frame}{Block, exampleblock, alertblock}
    \begin{block}{Block}
        This is a sample block.
    \end{block}

    \begin{exampleblock}{Example}
        This is a sample example block.
    \end{exampleblock}

    \begin{alertblock}{Alert}
        This is a sample alert block.
    \end{alertblock}
\end{frame}

\begin{frame}{Side-by-side blocks with overlays}
    \only<1->{%
        \begin{minipage}{0.48\textwidth}
            \begin{exampleblock}{Example 1}
                Content of example 1.
            \end{exampleblock}
        \end{minipage}}\hfill
    \only<2->{%
        \begin{minipage}{0.48\textwidth}
            \begin{exampleblock}{Example 2}
                Content of example 2.
            \end{exampleblock}
        \end{minipage}}
\end{frame}

\begin{frame}{Nested block}
    \begin{block}{Outer}
        Outer text.
        \begin{exampleblock}{Inner}
            Inner text, aligned with the outer block.
        \end{exampleblock}
    \end{block}
\end{frame}

%% ============================================================
\section{Columns \& Minipages}  [COL]
%% ============================================================
\begin{frame}{Columns: text $|$ figure}
    \begin{columns}[T]                       %% [T] aligns tops
        \begin{column}{0.48\textwidth}
            \begin{itemize}
                \item Left column text
                \item More text
                \item Even more
            \end{itemize}
        \end{column}
        \begin{column}{0.48\textwidth}
            \centering
            \includegraphics[width=0.8\linewidth]{figures/latexlogo.png}
            \par\smallskip
            {\footnotesize A caption under the figure.}
        \end{column}
    \end{columns}
\end{frame}

\begin{frame}{Minipages: three columns}
    \begin{minipage}[t]{0.31\textwidth}
        \begin{exampleblock}{A}
            \small First column.
        \end{exampleblock}
    \end{minipage}\hfill
    \begin{minipage}[t]{0.31\textwidth}
        \begin{exampleblock}{B}
            \small Second column.
        \end{exampleblock}
    \end{minipage}\hfill
    \begin{minipage}[t]{0.31\textwidth}
        \begin{exampleblock}{C}
            \small Third column.
        \end{exampleblock}
    \end{minipage}
\end{frame}

%% ============================================================
\section{Figures \& Tables}  [FIG] [TAB]
%% ============================================================
\begin{frame}{Single figure with caption and citation}
    \begin{figure}
        \centering
        \includegraphics[width=0.35\linewidth]{figures/latexlogo.png}
        \caption{\LaTeX\ logo}
        \label{fig:latexlogo}
        \citep{stallman1981emacs}
    \end{figure}
\end{frame}

\begin{frame}{Side-by-side figures}
    \begin{figure}
        \centering
        \begin{minipage}{0.45\linewidth}
            \centering
            \includegraphics[width=\linewidth]{figures/latexlogo.png}
            \caption{Left figure}
        \end{minipage}\hfill
        \begin{minipage}{0.45\linewidth}
            \centering
            \includegraphics[width=\linewidth]{figures/latexlogo.png}
            \caption{Right figure}
        \end{minipage}
    \end{figure}
\end{frame}

\begin{frame}{Table with booktabs and coloured header row}
    \centering
    \begin{tabular}{lrr}
        \toprule
        \iistheadrow Method & Speed (s) & Accuracy (\%) \\
        \midrule
        Baseline            & 12.4      & 88.1 \\
        Proposed            & 7.6       & 94.3 \\
        Oracle              & 5.1       & 99.0 \\
        \bottomrule
    \end{tabular}

    \medskip
    {\footnotesize Reference: \citep{stallman1985gnu}.}
\end{frame}

%% ============================================================
\section{Code}  [CODE]
%% ============================================================
%% Frames containing verbatim or lstlisting MUST use [fragile].
\begin{frame}[fragile]{Verbatin code}
\begin{verbatim}
def add(a, b):
    return a + b
\end{verbatim}
\end{frame}

\begin{frame}[fragile]{Listing with caption}
\begin{lstlisting}[language=Python, basicstyle=\ttfamily\small,
                   keywordstyle=\color{IISTPrimary}\bfseries,
                   commentstyle=\color{IISTMuted}\itshape,
                   stringstyle=\color{IISTExample},
                   numbers=left, numberstyle=\tiny,
                   frame=single, caption={Add two numbers}]
def add(a, b):
    # simple sum
    return a + b
\end{lstlisting}
\end{frame}

%% ============================================================
\section{Mathematics}  [MATH]
%% ============================================================
\begin{frame}{Inline and display math}
    Inline: $a = 1 + \theta$, and $\alpha^2 + \beta^2 = \gamma^2$.

    \medskip
    Display:
    \begin{equation}
        E = mc^2
        \label{eq:emc2}
    \end{equation}

    Refer back to Eq.~\eqref{eq:emc2}.
\end{frame}

\begin{frame}{Aligned equations, cases, and matrices}
    \begin{align}
        f(x)   &= x^2 + 2x + 1 \\
               &= (x+1)^2.
    \end{align}

    \begin{equation}
        g(x) =
        \begin{cases}
            x^2,   & x \ge 0, \\
            -x^2,  & x < 0.
        \end{cases}
    \end{equation}

    \begin{equation}
        A = \begin{pmatrix}
            a & b \\
            c & d
        \end{pmatrix},
        \qquad
        \det A = ad - bc.
    \end{equation}
\end{frame}

%% ============================================================
\section{Theorems}  [THM]
%% ============================================================
\begin{frame}{Theorem-like environments}
    \begin{assumption}
        The input sequence is i.i.d.\ and bounded.
    \end{assumption}

    \begin{proposition}
        Under the assumption above, the estimator is consistent.
    \end{proposition}

    \begin{remark}
        The assumption can be relaxed to stationarity.
    \end{remark}
\end{frame}

%% ============================================================
\section{Animations \& Overlays}  [ANIM]
%% ============================================================
\begin{frame}{Common overlay commands}
    \pause
    Line 1 appears first.

    \pause
    Line 2 appears on the next click.

    \pause
    Line 3 appears last.
\end{frame}

\begin{frame}{Only, uncover, visible, alt}
    \only<1>{Only on slide 1.}
    \only<2>{Only on slide 2.}

    \medskip
    \uncover<1->{Visible from slide 1, but takes space even when hidden.}

    \medskip
    \visible<2->{Visible from slide 2 onward, takes space throughout.}

    \medskip
    \alt<1>{Slide 1 version.}{Slide 2 and later version.}
\end{frame}

\begin{frame}{Onslide}
    \onslide<1->{Always shown.}
    \onslide<2->{Shown from slide 2.}
    \onslide<3>{Only on slide 3.}
\end{frame}

%% ============================================================
\section{Diagrams (TikZ)}  [TIKZ]
%% ============================================================
\begin{frame}{A simple TikZ diagram}
    \centering
    \begin{tikzpicture}[node distance=1.6cm, every node/.style={draw,
                        rounded corners, inner sep=4pt}]
        \node (a) {Input};
        \node (b) [right of=a] {Process};
        \node (c) [right of=b] {Output};
        \draw[->, thick, IISTPrimary] (a) -- (b);
        \draw[->, thick, IISTPrimary] (b) -- (c);
    \end{tikzpicture}
\end{frame}

%% ============================================================
\section{Algorithms}
%% ============================================================
\begin{frame}[label=algo]{Algorithm}
    \begin{algorithm}[H]
        \caption{Sum of $N$ numbers}
        $S = 0$\\
        \For{$i=1:N$}
        {
            $S = S + i$
        }
        Output $S$
    \end{algorithm}
\end{frame}

%% ============================================================
\section{Frames \& Options}  [FRAME]
%% ============================================================
\begin{frame}[plain]{Plain frame}
    No header, no footline. Good for title pages and closing slides.
\end{frame}

\begin{frame}[shrink=10]{Shrink to fit}
    Use \texttt{[shrink=<percent>]} when a frame overflows. Prefer
    editing the content over shrinking.
\end{frame}

\begin{frame}[label=reusable]{A labelled frame}
    This frame is labelled \texttt{reusable}.
\end{frame}

%% \againframe re-displays the frame above (useful in a summary).
\begin{frame}{Summary reuses the labelled frame}
    \againframe{reusable}
\end{frame}

\begin{frame}[allowframebreaks]{Long content that breaks across slides}
    \lipsum[1-6]
\end{frame}

%% ============================================================
\section{Citations \& Bibliography}  [CITE]
%% ============================================================
\begin{frame}{Footnotes and citations}
    \begin{block}{}
        ``GNU, which stands for Gnu's Not Unix, is the name for the
        complete Unix-compatible software system\footnote{%
            \citep{stallman1985gnu}.}''
    \end{block}
    EMACS is a real-time display editor that can be extended by the
    user while it is running.\footnote{\citep{stallman1981emacs}}
\end{frame}

\begin{frame}[allowframebreaks]{References}
    \bibliography{pres}
\end{frame}

%% ============================================================
%%  Backup slides  [BACKUP]
%% ============================================================
\appendix
\section*{Backup}
\begin{frame}{Backup: extra detail}
    Use \texttt{\textbackslash appendix} then plain frames for material
    that should not appear in the main flow but should be reachable
    during Q\&A.
\end{frame}

%% ============================================================
%%  Closing frames
%% ============================================================
\customframe{Questions?}{\emailid}
\customframe{Thank You!}{}

\end{document}
